%! Displaying and Analyzing Categorical Data — Unit 3, Lesson 3.1 — Warm-Up
% ONE PAGE, blank and key. Three items at 12pt. It also opens the unit (the first lesson of
% Unit 3, no separate launch lesson): one printed line says what the unit is about.
% THE NUMBERS ARE THE NOTES': item 1 gives 94, 54, 40 (row 1's bus bar, row 2's two bus bars);
% item 2 gives 47%, 45%, 50% (problems 8 and 9) and seeds the crux with bare fractions: the larger
% numerator is the smaller percent. Item 3 is row 1's one-variable data as a frequency table.
% The grade sizes are deliberately NOT printed with a context here: the hook stays open.
% Checked: 47% of 200 = 94; 45% of 120 = 54; 50% of 80 = 40; 94/200 = 47%; 54/120 = 45%;
% 40/80 = 50%; 26/200 = 13%.
\documentclass[12pt]{article}
\usepackage{saar-article}
\usepackage{saar-boxes}
\newcommand{\TallMath}[1]{$\displaystyle #1\rule[-1.4em]{0pt}{3.2em}$}

\begin{document}

\pageheader{Unit 3, Lesson 3.1}{Warm-Up}
% NAMESTRIP: no \namedateperiod here. The student writes their name once, on the
% cover this component is stapled behind.

\vspace{0.15in}

\boxguard
\begin{notesbox}{Unit 3 Starts Here --- Percents and a Table}
\normalsize
Unit~3 is about data sorted into groups and the chances those data give us. Today's tools are
ones you already own. Every number below comes back in the notes.

\begin{enumerate}[leftmargin=1.8em, itemsep=22pt, topsep=10pt]

  \item Find each amount.

        \vspace{4pt}
        \renewcommand{\arraystretch}{2.6}
        \begin{tabularx}{\linewidth}{@{} Y Y Y @{}}
        \toprule
        $47\%$ of $200$ & $45\%$ of $120$ & $50\%$ of $80$ \\
        \midrule
        \blank{2.6cm} & \blank{2.6cm} & \blank{2.6cm} \\
        \bottomrule
        \end{tabularx}

  \item Write each fraction as a percent.

        \vspace{4pt}
        \renewcommand{\arraystretch}{2.6}
        \begin{tabularx}{\linewidth}{@{} Y Y Y @{}}
        \toprule
        \TallMath{\dfrac{94}{200}} & \TallMath{\dfrac{54}{120}} & \TallMath{\dfrac{40}{80}} \\
        \midrule
        \blank{2.6cm} & \blank{2.6cm} & \blank{2.6cm} \\
        \bottomrule
        \end{tabularx}

        \vspace{8pt}
        Look at $\dfrac{54}{120}$ and $\dfrac{40}{80}$. Which one has the larger top number? Which
        one is the larger percent?
        \par\writelines{2}

  \item Oakdale High asked $200$ students how they get to school.

        \vspace{4pt}
        \renewcommand{\arraystretch}{1.8}
        \begin{tabularx}{\linewidth}{@{} X Y Y Y Y @{}}
        \toprule
        \textbf{How they get there} & \textbf{Bus} & \textbf{Car} & \textbf{Walk} & \textbf{Total} \\
        \midrule
        \textbf{Students} & $94$ & $80$ & $26$ & $200$ \\
        \bottomrule
        \end{tabularx}

        \vspace{10pt}
        How many students come by car? \blank{2.6cm}

        \vspace{12pt}
        What percent of the $200$ students walk? \blank{2.6cm}

\end{enumerate}
\end{notesbox}

\end{document}
